% !TeX root = ../papers/nc-p3.tex
\section{三重色と二重色が隣接する場合の全排除}
\label{sec:klein-mixed-positive}
\begin{theorem}\label{thm:klein-mixed-positive}
正則表現型四元群窓において、二つの隣接切断の色階数がそれぞれ
$(3,1,1,1)$ と $(2,2,1,1)$ の置換である標準14641型Koszul左右AS正則代数は存在しない。
\end{theorem}
\begin{proof}[台と五つの長方形]
文字を平行移動し共通自己同型を施し、$M_0$ が階数3、$N_0,N_1$ が階数2とする。
$M_1,M_2,M_3$ の非空長方形を選ぶと各 $N_e$ の12セルが決まる。残る4セルの全ての補完を調べ、
階数1では長方形、階数2では零積で定まる行・列直線の同値閉包を使う。
最後に $M_0$ に前節のマッチング検査を施す。台の必要条件を通る数は109034088、
四層の自由辺の指数増大検査を通る数は15000である。
うち256個では五長方形の56係数が全て非零である。残る14744個を
片側欠色窓と同じ文字置換・平行移動・巡回・反転で集めると312軌道となる。
階数2の行射影には直線同値閉包、階数3にはマッチング上界を使い、
連続純関係と一文字商で検査すると、前節と同じビット順で
$\mathtt{0ff00f0f00ffcfff}$ と $\mathtt{0fff0fff0fffcfff}$ の二代表だけが残る。
\end{proof}

\begin{proof}[五長方形が全て四行四列の場合]
階数1の指数行列の非零Smith対角は42個の1であり、対角基底変更は階数13である。
従って56係数の商には一変数 $a\ne0$ が残る。未定の8成分を $x_0,\ldots,x_7$ とすれば、
Smith変換の逆で得る三行列 $N_0,N_1,M_0$ は順に
\[
\begin{pmatrix}x_0&a&a&1\\1/a&x_3&1/a&1\\1/a&a&x_4&1\\1&1&1&x_7\end{pmatrix},\quad
\begin{pmatrix}a&x_1&1&a\\x_2&1/a&1&1/a\\a&1/a&1&x_5\\1&1&x_6&1\end{pmatrix},\quad
\begin{pmatrix}x_0&x_1&1&1\\x_2&x_3&1&1\\1&1&x_4&x_5\\1&1&x_6&x_7\end{pmatrix}.
\]
それぞれの $N$ の四変数を $(u,v,f,h)=(x_0,x_3,x_4,x_7)$ または
$(x_1,x_2,x_5,x_6)$ とする。$a(a^2-1)$ を可逆にすると、三次小行列式のイデアルは
$(u-v,uh-1,f-h,(h-a)(ah-1))$ に等しい。検算は可逆化変数を加えて両側の
Gröbner基底を比較するので、$a^2\ne1$ の特殊値も含む。よって $x_7,x_6$ はそれぞれ
$a,1/a$ の一方である。両者が等しければ $M_0$ は階数1、異なれば行列式 $-(a-1/a)^4$ で階数4になる。

$a=\pm1$ の各々では上記小行列式のイデアルは、異なる二変数 $y,z$ の全ての積
$(y-a)(z-a)$ で生成される。全変数が $a$ のときそれぞれの $N$ と $M_0$ は階数1である。
各 $N$ が階数2ならそれぞれ一セルだけが変わり、$N_0$ 側の変更は対角のある位置 $i$ にある。
$M_0$ が階数3なら二つの変更セルの列ベクトルと元の列直線は一次独立で、その像は $e_i$ を含む。
$N_0$ の像も $e_i$ を含むので、色零の純平方が連続して標準型に反する。
\end{proof}

\begin{proof}[二つの疎な代表]
五長方形の非零行・列の先頭を基準として同じSmith計算を行う。第一代表では30係数に対し
非零対角17個が全て1で自由変数は $a,b$、第二代表では42係数に対し28個が全て1で自由変数は $a$ である。
残る八つの係数 $x_i$ は語の辞書順に取り、ここでは全て非零である。
第一代表では二つの $N$ の三次小行列式から
\[
 x_1=x_2=x_5=1,\quad ax_1=b,\quad x_0=b,\quad ax_3=1,\quad
 bx_4=a,\quad a^3x_0=1
\]
を得る。従って $a=b$, $a^4=1$。第二代表では同じ検査で
$(x_0,x_1,x_2,x_3,x_4,x_5)=(a,a,1,1,a,a)$ かつ $ax_1=1$ となる。
いずれも $a=\pm1$ では $M_0$ の階数が高々2である。

第一代表の $a^2=-1$ では $\det M_0=-4(x_6+x_7)$ なので $x_7=-x_6$ に限られる。
このとき始点1、長さ4、文字次数2の商次元は $k(x_6)$ 上で9である。
実際、総文字次数2の64語を列とし、三位置の二次縮約の全列に外側の二文字を付けて
作る関係行列の階数は55である。特殊化によって階数は増えない。
一方、標準Koszul解消の文字別Euler式はこの次数で8を要求する。
具体的には長さ4で正則表現の畳み込みから来る寄与が相殺し、
長さ2の二関係の寄与 $\sum_b r_b r_{2+b}=8$ が残る（$r=(2,2,1,1)$）。
最上項は総文字次数零にだけ寄与する。従ってこの最後の候補も全点で排除される。

\path{computations/verify_klein_mixed_positive.py} と付属のC++17列挙が、
全台の数、二つの整数正規化、三次小行列式と例外値、最後の文字別階数を再現する。
\end{proof}
